Practica de trivia con API (simulada).
Objetivos: Desarrollar una aplicación web utilizando HTML, CSS, JavaScript (sin framework) que permita jugar a una trivia de preguntas y respuestas. 
Consumiendo datos desde un archivo JSON local (simulando una API).

1- Recursos provistos:
Se entrega un archivo: 
    •	Trivia_240_preguntas.json
Ese archivo contiene:
    •	6 categorías.
    •	40 preguntas por categoría.
Cada pregunta tiene la siguiente estructura:

{
"id": "HIS_001",
"pregunta": "Texto de la pregunta",
"correcta": "Respuesta correcta",
"incorrectas": ["Opción 1", "Opción 2", "Opción 3"],
"nivel": "fácil"
}

2- Requisitos funcionales:
1.	Seleccionar categoría:
    •	Al iniciar la aplicación, el usuario debe poder seleccionar una categoría.
    •	No se puede comenzar el juego sin seleccionar una categoría.

2.	Carga de datos:
    •	Las preguntas deben cargarse utilizando fetch() desde el archivo JSON local.
    •	Si ocurre un error al cargar los datos:

1-	Mostrar un mensaje en pantalla.

2-	No permitir continuar.

3.	Selección de preguntas:
•	Una vez elegida la categoría:
    Seleccionar 5 preguntas aleatorias, no deben repetirse.

4.	Desarrollo del juego:
•	Por cada pregunta:
    1.	Mostrar: 
             # Texto de la pregunta.
             # 4 opciones (1 corecta + 3 incorrectas).

    2.	Las respuestas deben:
             # Estar mezcladas aleatoriamente.
             # No mantener siempre la misma posición.


5.	Tiempo por pregunta:
    •	Cada pregunta tiene un máximo de 5 segundos.
    •	Debe mostrar un contador visible en pantalla.
   Comportamiento obligatorio:
            •	Si el usuario responde antes: (Se evalúa la respuesta; Se para a la siguiente pregunta)
            •	Si el tiempo llega a 0: (La pregunta se considera no respondida; Se pasa automáticamente a la siguiente)

6.	Restricciones de interacción:
•	Una vez seleccionada una respuesta: (No se puede cambiar; Se debe deshabilitar las demás opciones).

7.	Resultado final: Al finalizar las 5 preguntas, se debe mostrar:
•	Total de preguntas: 5
•	Cantidad de respuestas: (Correctas, Incorrectas, No respondidas)

3-Requisitos de interfaz:
•	Interfaz clara y ordenada.
•	Debe diferenciar visualmente: 
        a.	Respuesta correcta.
        b.	Respuesta incorrecta.
•	Mostrar el progreso (Ej: “Pregunta 2 de 5”).
•	Mostrar el tiempo restante en cada pregunta.

4-Requisitos técnicos:
•	Usar únicamente: HTML, CSS, JavaScript (vainilla).
•	No usar librerías externas.
•	Uso obligatorio de:
        a.	Fetch ()
        b.	setTimeout () y/o setInterval ().
•	Manipulación dinámica del DOM

5-Restricciones:
•	No hardcodear preguntas en el código.
•	No repetir preguntas.
•	No permitir múltiples respuestas por pregunta.

6-Resultado esperado:
Una aplicación funcional donde el usuario:
•	Selecciona una categoría
•	Responde 5 preguntas con límite de tiempo
•	Visualiza su resultado final

7-Puntos adicionales (no obligatorios):
•	Botón para reiniciar el juego.
•	Animaciones (CSS).
•	Selección de dificultad (nivel).
•	Guardar puntaje en localStorage